% =============================================================================
% Manuscript: Spectral Co-Design of Agrivoltaic Modules for Multifunctional Net Energy Benefit:
%             A Quantum-Guided Digital Twin
% Target: Applied Energy (Elsevier)
% Date:    September 27, 2026 (Applied Energy submission package)
% Authors: Teguia Kouam S.C., Goumai Vedekoi T., Tchapet Njafa J.-P.,
%          Nguenang J.-P., Nana Engo S.G.
% Note: submission package copied from Submission_Package_Nature_Energy_Manuscript;
%       base files preserved as *_BASE.pdf. Converted to elsarticle (preprint).

\documentclass[preprint,12pt]{elsarticle}
\usepackage[utf8]{inputenc}
\usepackage[T1]{fontenc}
\usepackage{amsmath,amssymb,amsfonts}
\usepackage{graphicx}
\graphicspath{{figures/}}
\usepackage{booktabs}
\usepackage{siunitx}
\usepackage{physics}
% mhchem v4 for chemical formulas (\ce{Pb^{2+}}, \ce{CO2}, etc.)
\usepackage[version=4]{mhchem}
% siunitx v3 ↔ physics compatibility: physics defines \qty, siunitx v3 redefines it.
% This shim ensures siunitx's \qty takes precedence at begin-document time.
\AtBeginDocument{\RenewCommandCopy\qty\SI}
\usepackage{xcolor}
\usepackage[colorlinks=true,linkcolor=blue,citecolor=blue,urlcolor=blue]{hyperref}
\usepackage[capitalise,nameinlink]{cleveref}
\usepackage{xr}
\externaldocument{AppliedEnergy_SM_2609}
% (elsarticle manages page geometry; the geometry package is intentionally not loaded)
\usepackage{lineno}

% ---- siunitx v3 configuration ----
\sisetup{
    locale = US,
    detect-all = true,
    group-minimum-digits = 4,
    group-digits = integer,
    separate-uncertainty = true,
    multi-part-units = single,
    range-phrase = {--},
    range-units = single,
    number-unit-product = \,,
    allow-number-unit-breaks = true,
}
\DeclareSIUnit{\yr}{yr}
\DeclareSIUnit{\USD}{USD}
\DeclareSIUnit{\tonne}{t}
\DeclareSIUnit{\ha}{ha}
\DeclareSIUnit{\hectare}{ha}
\DeclareSIUnit{\MWh}{MWh}
\DeclareSIUnit{\kWh}{kWh}
\DeclareSIUnit{\cubed}{^{3}}
\DeclareSIUnit{\m}{\meter}

% Line numbering disabled for production length; re-enable with \linenumbers for review.
\emergencystretch=2.5em

% eqphift removed with detailed equations (now in SM)

\begin{document}
\linenumbers

\begin{frontmatter}

\title{Spectral Co-Design of Agrivoltaic Modules for Multifunctional Net Energy Benefit: A Quantum-Guided Digital Twin}

\author[udla]{Steve Cabrel Teguia Kouam\corref{cor1}}
\author[uy1]{Theodore Goumai Vedekoi}
\author[uy1]{Jean-Pierre Tchapet Njafa}
\author[udla]{Jean-Pierre Nguenang}
\author[uy1]{Serge Guy Nana Engo}

\cortext[cor1]{Corresponding author. Email: steve.teguia@facsciences-uy1.cm}

\address[udla]{Department of Physics, Faculty of Science, University of Douala, Cameroon}
\address[uy1]{Department of Physics, Faculty of Science, University of Yaound\'{e} I, Cameroon}
% ORCID identifiers to be attached by authors before upload.

\begin{abstract}
Spectral co-design offers a route beyond the system-level trade-off of static agrivoltaics, in which semi-transparent modules resolve the food--energy land conflict only by uniformly sacrificing electrical output or crop productivity. Here we present a quantum-guided agrivoltaic stack in which a semi-transparent organic photovoltaic layer couples to a plasmonic nanoparticle-on-mirror (NPoM) cavity that actively shapes the transmitted spectrum. By transmitting narrow bands at \qtylist{750;820}{\nm}, the system prepares polaron-dressed states in the Fenna--Matthews--Olson photosynthetic complex and extends exciton coherence by \qty{50}{\percent}; confining the plasmonic sentinel layer to \qty{1}{\percent} of the area then preserves \qty{97.1}{\percent} of the device-level excitation-transfer yield (a device metric, not crop yield) at a \qty{91.8}{\percent} local transport penalty on the sentinel. The NPoM sentinel layer enables multiplexed SERS and quantum dot diagnostics of agricultural contaminants and crop stress. On the macroscopic scale, the shield reduces crop evapotranspiration by \qty{28}{\percent} and yields a Net Ecological Benefit of \qty{72.6}{\kg\of{\ce{CO2e}}\per\m\squared\per\yr}. A BB84 quantum link secures the IoT telemetry. Finally, techno-economic analysis of a Cameroonian micro-module gives a payback period of \qty{4.32}{\yr}, identifying conditions under which the multifunctional net benefit stays positive for the modeled \qty{500}{\m\squared} cooperative case.
\end{abstract}

\begin{keyword}
agrivoltaic spectral management \sep net ecological benefit \sep FAO-56 microclimate \sep plasmonic nanoparticle-on-mirror cavity \sep SERS diagnostics \sep techno-economic analysis \sep quantum-guided digital twin
\end{keyword}

%Graphical abstract
% \begin{graphicalabstract}
% \includegraphics[width=\linewidth]{Graphical_Abstract_wide.png}
% \end{graphicalabstract}

\end{frontmatter}
% Frontmatter spacing handled by elsarticle; no extra page enlarge needed without lineno.

% =============================================================================
% INTRODUCTION
% =============================================================================

\section{Introduction}

Agrivoltaics co-locates photovoltaic (PV) systems and agricultural cultivation to mitigate the land-use competition between food production and renewable energy generation~\cite{Dupraz2011}. Land-use optimization and mobile-PV foundations for this field were established in \textit{Applied Energy}~\cite{Valle2017,Amaducci2018}. Recent syntheses systematize overhead, vertical, and tracking typologies with land-equivalent-ratio and levelized-cost metrics~\cite{Mahim2025}.
Across these studies, three gaps recur: (i)~the optical architecture is static, so spectral light-sharing is not co-optimized with crop photophysiology~\cite{Thompson2020,Zhao2023}; (ii)~techno-economic and lifecycle accounts remain decoupled from the underlying transport physics~\cite{Ravilla2024,Bellone2024}; and (iii)~the PV layer is treated purely as an energy harvester, with crop-stress and contaminant sensing bolted on as separate hardware~\cite{CiallaMay2024,Pook2025,KaurJeet2026}. Dynamic, physics-guided spectral management and multifunctional co-design have not yet been transposed into a single operating agrivoltaic stack.

The photosynthetic apparatus of green plants and bacteria has evolved under a spectrally structured photon flux, where specific wavelengths drive charge separation while others dissipate as heat~\cite{Blankenship2011}.
This spectral specificity lets the PV layer transmit only the photosynthetically active wavelengths while absorbing the remainder, converting the food-energy conflict into a problem of spectral co-design.

The Fenna--Matthews--Olson (FMO) complex of green sulfur bacteria serves as the model system for spectral light-sharing~\cite{Fenna1974}.
Its seven bacteriochlorophyll $a$ sites have been extensively characterized through two-dimensional electronic spectroscopy, revealing oscillatory coherence signatures that persist for hundreds of femtoseconds at cryogenic~\cite{Engel2007} and room temperature~\cite{Panitchayangkoon2010}.
Although the functional relevance of these quantum beats remains debated, consensus has emerged that vibronic coupling --- the hybridization of electronic transitions with underdamped protein vibrations --- governs the flow of excitation energy~\cite{Chin2013,Christensson2012}.
Non-Markovian frameworks such as the hierarchical equations of motion~\cite{Tanimura1989} and the process tensor hierarchy of pure states (PT-HOPS)~\cite{Varvelo2021,Citty2024} capture these effects with demonstrated numerical convergence. Device-adjacent uses of exact open-quantum-system methods remain rare: the closest precedent is a general HEOM library applied to light-harvesting and control problems without a photovoltaic product~\cite{Lambert2023}. Photonic-quantum solar design and vibronic-screening workflows likewise stop short of an operating device~\cite{LiuFang2026,Motlagh2024}.

Previous work demonstrated that spectral filtering of the excitation pulse --- aligning its spectrum with underdamped vibronic resonances at \qty{750}{\nm} and \qty{820}{\nm} --- extends coherence lifetimes by \qty{20}{\percent} to \qty{50}{\percent} in the 7-site FMO complex~\cite{TeguiaKouam2026}.
This selective vibronic excitation prepares polaron-dressed states that decouple from the dissipative protein-solvent bath. That companion paper, published in the \textit{Journal of Physical Chemistry Letters}, reports the analysis of the filter, the excitation-transport rate, and the single-pass conversion efficiency~\cite{TeguiaKouam2026}; the present work adds the 8-site trapping model, the monolithic OPV--NPoM hardware concept, the SERS field calibration, and the Net Ecological Benefit and lifecycle assessment developed below.
Single-photon experiments have confirmed that quantum coherence persists in individual light-harvesting complexes under ambient conditions~\cite{Li2023}, and quantum phase synchronization via exciton-vibrational energy dissipation has been identified as a mechanism sustaining long-lived coherence~\cite{Zhu2024}.
Here we extend this principle in three directions.

First, we introduce an 8-site FMO model with a non-Hermitian reaction center trap on sites 3 and 4, enabling quantitative computation of the forward transfer yield $\Phi_{\mathrm{FT}}$ as the time-integrated population captured by the reaction center.
Second, we embed the FMO complex in a nanoparticle-on-mirror (NPoM) plasmonic cavity~\cite{Chikkaraddy2016} that couples the biological light-harvesting system to the organic photovoltaic (OPV) layer through strong coherent coupling, forming a $9\times 9$ dressed Hamiltonian.
Third, we specify dynamic photo-protection mechanisms --- Floquet Stark detuning and flux-dependent optical limiting --- as a control law that mimics the non-photochemical quenching used by plants under high light stress.

We map this microscopic optimization onto a field-scale lifecycle assessment.
We parameterize crop evapotranspiration using the FAO-56 Penman--Monteith model~\cite{Allen1998} under the spectrally selective OPV shield, compute the Net Ecological Benefit (NEB) across three scenarios, and evaluate socioeconomic access through cooperative ownership models and BB84 quantum key distribution~\cite{Bennett1984} for secure IoT telemetry.
This multi-scale perspective~\cite{Vandamme2025} links angstrom-scale vibronic couplings to square-meter canopy yields, providing a unified modeling framework for quantum-guided agrivoltaics (\cref{fig:system_setup}).

The central question of this work is falsifiable: \emph{Can spectral co-design guided by non-Markovian quantum dynamics deliver a positive multifunctional net energy benefit (power $+$ water $+$ carbon $+$ diagnostics) in a field-constrained agrivoltaic stack, and which barriers remain physical rather than economic?}
Three subordinate questions structure the Results. (Q1) Device asks what the NPoM co-design costs and where the cost lands, quantified by the mode-volume scan (\Cref{tab:npom_volume_scan}) and located by the canopy-integration analysis. (Q2) Physics asks which microscopic states carry the trade-off: exact PT-HOPS dynamics localize the loss channels, and Floquet--Stark detuning with flux-dependent optical limiting specifies a control law under which they could be managed dynamically (its yield-level validation, $\Phi_{\mathrm{FT}}$ with the switch engaged vs.\ disengaged, is an Outlook step). (Q3) System asks whether the stack pays for itself, once coupled FAO-56 microclimate, cooperative techno-economics, and BB84-secured telemetry internalize every function into a single Net Ecological Benefit account.
Three falsifiable claims anchor the answer. The first, (F1) Microscopic, is that dual-band spectral filtering at \qty{750}{\nm} and \qty{820}{\nm} prepares polaron-dressed states that extend exciton coherence by \qty{50}{\percent} and raise the forward transfer yield by \qty{25}{\percent} at \qty{295}{\kelvin} --- both established in the companion 7-site study~\cite{TeguiaKouam2026} and carried over here under the weak site-8 coupling assumption.
The second, (F2) Device, is that confining the NPoM cavity to a \qty{1}{\percent} sentinel fraction preserves \qty{97.1}{\percent} of the global canopy trapping yield ($\Phi_{\mathrm{FT}}^{\mathrm{global}} = 0.971$) while retaining single-complex SERS diagnostics, despite the $>\qty{90}{\percent}$ local transport penalty quantified across the mode-volume scan.
The third, (F3) System, is that the FAO-56 smart shield cuts evapotranspiration by \qty{28}{\percent} ($\pm\qty{10}{\percent}$ microclimate uncertainty) and, combined with OPV generation and cooperative economics, yields a Net Ecological Benefit of \qty{72.6}{\kg\of{\ce{CO2e}}\per\m\squared\per\yr} and a \qty{4.32}{\yr} payback at \qty{30}{\percent} subsidy for the modeled \qty{500}{\m\squared} cooperative case.
We answer affirmatively but conditionally: the bounds (single-exciton treatment, site-independent baths, modeled fouling and matrix penalties, $\pm\qty{10}{\percent}$ microclimate band) are stated explicitly in the Limitations section of the Discussion.
The headline outcome, pictured in \Cref{fig:flat_canopy}, compresses this answer to one sentence: a \qty{91.8}{\percent} local transport penalty confined to a \qty{1}{\percent} sentinel area costs 0.9 points of canopy-scale trapping yield (\num{0.980} to \num{0.971}) while adding single-complex SERS diagnostics to the same footprint.
To keep a stack of this breadth auditable, each layer --- quantum dynamics, microclimate, economics, and security --- is validated separately in the Supplementary Material (\Cref{SI-sec:convergence}--\Cref{SI-sec:roadmap}), and every headline number below is traceable to a single figure or table.

The value chain of the quantum layer is worth stating explicitly, because the energy-systems reader is entitled to ask what the quantum hardware buys. It contributes three concrete functions: (i) the dual-band passband co-design that the FAO-56 and lifecycle accounts price into the module's transmittance budget; (ii) multiplexed SERS and quantum-dot diagnostics delivered on the \qty{1}{\percent} sentinel footprint at no additional canopy area; and (iii) preservation of \num{0.971} of the device-level excitation-transfer yield despite the local plasmonic penalty. On the cost side, the NPoM sentinel and the BB84 link appear as capital and operating line items in the techno-economic account. The transport gain itself ($+\qty{25}{\percent}$ forward transfer yield) is deliberately not monetized: no coherence-derived revenue enters the NEB or payback chains.

% =============================================================================
% RESULTS
% =============================================================================

\section{Results}

\subsection{Quantum-enabled spectral co-design and NPoM sentinel}

The spectral co-design rests on an 8-site FMO model with a non-Hermitian reaction-center trap on sites~3 and~4 (full Hamiltonian, site energies, and 28-pair coupling matrix in \Cref{SI-sec:hamiltonian}). Dual-band transmission at \qtylist{750;820}{\nm} prepares polaron-dressed states that extend the $l_1$-norm coherence lifetime from \qty{280(25)}{\fs} (broadband) to \qty{420(35)}{\fs} (filtered) --- a \qty{50}{\percent} increase --- and raise the forward transfer yield $\Phi_{\mathrm{FT}}$ from \num{0.71(2)} to \num{0.89(3)} (\Cref{fig:quantum_dynamics}; PT-HOPS/SBD propagation details and convergence in \Cref{SI-sec:convergence}).

The OPV--FMO interface is mediated by a nanoparticle-on-mirror (NPoM) cavity that forms a $9\times 9$ dressed Hamiltonian. Mode-volume scans (\Cref{tab:npom_volume_scan} and \Cref{SI-sec:comparative_hdf5}) show that strong plasmon--exciton coupling suppresses local transport to $\Phi_{\mathrm{FT}}^{\mathrm{NPoM}} \approx 0.080$ at $V_{\mathrm{mode}} = \qty{1.2}{\nm\cubed}$ while enabling relative Purcell-scaled SERS enhancement factors of \num{100} at $V_{\mathrm{mode}} = \qty{0.8}{\nm\cubed}$ down to \num{44} at $V_{\mathrm{mode}} = \qty{1.2}{\nm\cubed}$ (a factor-\num{48} enhancement range across the full seven-point scan, from \num{33} to \num{1.6e3}; tabulated endpoints here, full scan in \Cref{SI-sec:comparative_hdf5}), consistent with absolute enhancements $\geq 10^6$ reported for sub-nm NPoM gaps. The integrated transmitted power in the two \qty{100}{\per\cm} passbands is $\sim\qty{10.2}{\W\per\m\squared}$ ($\sim\qty{1.0}{\percent}$ of AM1.5G), priced into the module transmittance budget. The \qty{91.8}{\percent} local suppression relative to the passive baseline ($\Phi_{\mathrm{FT}}^{\mathrm{passive}} = 0.980$) is the price of the diagnostic channel; it is confined to a sentinel area fraction $\alpha = 0.01$.

\begin{table}[htbp]
\centering
\caption{\textbf{NPoM mode-volume sensitivity of forward transfer yield (selected points).} Passive baseline $\Phi_{\mathrm{FT}} = 0.980$ (assumed capture-efficiency cap, a model input rather than a trajectory integral). Full seven-point scan, SERS enhancement factors, and $n$-trajectory provenance in \Cref{SI-sec:comparative_hdf5} ($n=100$ at $V_{\mathrm{mode}}=\qty{0.8}{\nm\cubed}$, $n=20$ at $V_{\mathrm{mode}}=\qty{1.2}{\nm\cubed}$, $n=2$ otherwise). Production parameters: $L=8$, $K=2$, $\Delta t = \qty{0.2}{\fs}$, $T=\qty{295}{\kelvin}$.}
\label{tab:npom_volume_scan}
\begin{tabular}{S[table-format=1.1] S[table-format=1.4] S[table-format=-2.1]}
\toprule
{$V_{\mathrm{mode}}$ (\unit{\nm\cubed})} & {$\Phi_{\mathrm{FT}}$} & {$\Delta$ vs.\ baseline (\unit{\percent})} \\
\midrule
0.2 & 0.0505 & -94.8 \\
0.8 & 0.0768 & -92.2 \\
1.2 & 0.0799 & -91.8 \\
1.4 & 0.0791 & -91.9 \\
\bottomrule
\end{tabular}
\end{table}

\subsection{Global canopy yield integration}

Because the NPoM SERS diagnostic mode suppresses local excitonic transport through plasmon population loading, it need only cover a \qty{1}{\percent} sentinel fraction $\alpha$ of the canopy area. The remaining \qty{99}{\percent} operates in passive OPV mode, preserving the near-baseline forward transfer yield $\Phi_{\mathrm{FT}}^{\mathrm{passive}} = 0.980$. The area-weighted \emph{global} canopy trapping yield is therefore:
\begin{equation}
  \Phi_{\mathrm{FT}}^{\mathrm{global}} = \alpha\,\Phi_{\mathrm{FT}}^{\mathrm{NPoM}} + (1-\alpha)\,\Phi_{\mathrm{FT}}^{\mathrm{passive}},
  \label{eq:phi_global}
\end{equation}
yielding $\Phi_{\mathrm{FT}}^{\mathrm{global}} = 0.971$ at $\alpha = 0.01$ and $\Phi_{\mathrm{FT}}^{\mathrm{NPoM}} = 0.080$ (or $0.0799$ precisely); device-level excitation transfer is thus maintained near baseline even though transport is suppressed locally on the sentinel.
$\Phi_{\mathrm{FT}}^{\mathrm{global}}$ is a device-level metric: crop productivity is governed by the photosynthetically active radiation transmitted through the module (the AVT/PAR budget), not by excitonic trapping in the FMO layer --- the two are optically coupled only through the shared passband design, which is why no coherence-derived gain is booked in the crop revenue.
Because \Cref{eq:phi_global} is linear in $\alpha$, the flatness of \Cref{fig:flat_canopy}b is an accounting identity rather than an independent dynamical result; what carries physical content are the two reference endpoints ($\num{0.080}$, measured on the sentinel run, and $\num{0.980}$, the passive capture-efficiency cap) that enter it.
The identity converts those endpoints into a design envelope: solving \Cref{eq:phi_global} for $\alpha$ gives a sentinel fraction of up to \qty{3.3}{\percent} before the global yield falls below $0.95$, and up to \qty{8.9}{\percent} before it falls below $0.90$.
The diagnostic area therefore has a factor $\sim 3$ of headroom over the \qty{1}{\percent} assumed here --- the co-design lever is not a better mode volume but how much area one is willing to spend on sensing.

\subsection{Dynamic photo-protection}
Under high irradiance ($\gtrsim\qty{850}{\W\per\m\squared}$) a Floquet--Stark detuning of the OPV--FMO coupling, combined with flux-dependent optical limiting of the dual-band passbands, provides a plant-inspired non-photochemical-quenching analogue (\Cref{fig:npq_switch}). The control law is specified in \Cref{SI-sec:npq}; the production trajectories reported above are propagated without the drive engaged. The design therefore supplies a switchable photo-protection channel whose yield-level validation remains an Outlook step.

\subsection{In situ SERS diagnostics and environmental robustness}
The same NPoM sentinel that carries the strong-coupling penalty enables multiplexed surface-enhanced Raman spectroscopy of pesticides, heavy metals, and crop-stress markers at relative Purcell-scaled enhancement factors (\num{100} at $V_{\mathrm{mode}} = \qty{0.8}{\nm\cubed}$, \num{44} at $V_{\mathrm{mode}} = \qty{1.2}{\nm\cubed}$). Soiling and matrix penalties are parameterized by a DynamicCalibrator module (soiling \qty{0.005}{\per\day}, floor \num{0.75}; $\times 10$ LOD field penalty on raw matrices); the resulting diagnostic availability remains above \qty{90}{\percent} under the modeled tropical fouling schedule (\Cref{SI-sec:comparative_hdf5}). No additional canopy area is required beyond the $\alpha=0.01$ sentinel fraction already accounted for in $\Phi_{\mathrm{FT}}^{\mathrm{global}}$.

\subsection{Smart shield microclimate: FAO-56 evapotranspiration}
Under the dual-band OPV+NPoM transmittance the FAO-56 Penman--Monteith calculation yields a \qty{28}{\percent} reduction in crop evapotranspiration relative to the open-field baseline (central value; $\pm\qty{10}{\percent}$ microclimate band). For the \qty{500}{\m\squared} module this corresponds to a water credit of \qty{460}{\L\per\m\squared\per\yr}. Full meteorological forcing, albedo, and aerodynamic resistance parameters are tabulated in the Supplementary Material.

\subsection{Socioeconomic viability and cooperative payback}
Techno-economic analysis of the Cameroonian cooperative case produces a payback of \qty{4.32}{\yr} at a \qty{30}{\percent} blended-finance subsidy (unsubsidized \qty{6.17}{\yr}; local-material scenario \qty{2.68}{\yr}). Net Ecological Benefit reaches \qty{72.6}{\kg\of{\ce{CO2e}}\per\m\squared\per\yr} (carbon revenue \qty{363}{\USD\per\yr}, $\sim\qty{1.4}{\percent}$ of cash flow). Capital, OPEX, revenue streams, and sensitivity tables are given in \Cref{SI-sec:socioeconomic}; the quantum layer enters the cash-flow only as a cost item (sentinel + BB84) and as preserved device-level function.

\subsection{Secure IoT telemetry via BB84}
A prepare-and-measure BB84 link secures the high-stakes irrigation and crop-protection commands. At a realistic channel noise of \qty{4.8}{\percent} the protocol generates a \qty{256}{\bit} key per session; the abort decision uses the \qty{95}{\percent} Wilson upper confidence bound on a $\geq 128$-bit sample against the \qty{11}{\percent} Shor--Preskill threshold. On abort, a fail-safe local irrigation routine limited to \qty{5}{\L\per\m\squared\per\day} is engaged (\Cref{SI-sec:qkd}).

\subsection{Experimental testability}
Key observables---coherence lifetime extension, dual-band $\Phi_{\mathrm{FT}}$ gain, SERS enhancement, and ET reduction---are accessible with existing 2DES, broadband transient absorption, portable Raman, and lysimeter instrumentation. A staged validation roadmap (lab $\to$ controlled greenhouse $\to$ cooperative field trial) is outlined in \Cref{SI-sec:roadmap}.

\section{Discussion}

Our results show that spectral co-design guided by exact non-Markovian dynamics can deliver a modeled multifunctional net benefit (power + water + carbon + diagnostics) under the constraints of a \qty{500}{\m\squared} cooperative micro-module.

\subsection{Mechanism generality and prior art}
The dual-band resonance condition $\Delta E_{nm} \approx \hbar\omega_{\mathrm{vib}} \pm J_{nm}$ is not FMO-specific (\Cref{SI-sec:resonance}). A minimal three-site trimer recovers the same qualitative yield gain (\Cref{SI-sec:3site}), indicating that vibronically coupled aggregates (synthetic J-aggregates, plant LHCs, or conjugated polymer domains) can benefit from the same spectral protection. Relative to prior agrivoltaic optical and TEA studies~\cite{Valle2017,Amaducci2018,Mahim2025,Ravilla2024,Bellone2024}, the present stack is distinguished by (i)~an explicitly co-optimized dual-band passband, (ii)~a quantified local-to-global transport trade-off (\num{0.980} $\to$ \num{0.971} at $\alpha=0.01$), and (iii)~integration of SERS diagnostics and BB84 telemetry into the same area and cost accounts. A fuller comparative positioning against recent agrivoltaic TEA and optical studies is provided in the Supplementary Material.

\subsection{Limitations}
Four bounds are material. (1)~The single-exciton, site-independent-bath treatment omits multi-exciton annihilation and site-specific spectral densities. (2)~The NPoM transport penalty exceeds \qty{90}{\percent} locally; the global yield of \num{0.971} is the value at the \qty{1}{\percent} sentinel fraction (headroom to \qty{3.3}{\percent} before $0.95$, \qty{8.9}{\percent} before $0.90$). (3)~Microclimate and ET reductions carry a $\pm\qty{10}{\percent}$ uncertainty band from the FAO-56 parameterization under the modeled shield transmittance. (4)~Fouling, matrix-penalty, and encapsulation lifetime are modeled rather than field-validated; the DynamicCalibrator and soiling factors are therefore provisional. No coherence-derived crop-yield or revenue term enters the NEB or payback calculation; the quantum layer appears only as a cost (sentinel + QKD) and as preserved device-level function.

\subsection{Outlook}
Priority next steps are two-dimensional electronic spectroscopy validation of the dual-band protocol, low-loss dielectric antenna alternatives that reduce the local transport penalty, field fouling trials of the \qty{1}{\percent} sentinel network, and a localized Au/graphene supply chain for Sub-Saharan manufacturing. With local materials the unsubsidized capital cost falls from \qty{322}{\USD\per\m\squared} to \qty{140}{\USD\per\m\squared} (payback \qty{2.68}{\yr}), shifting the dominant barrier from economics to process engineering.

\section{Conclusions}

We asked whether spectral co-design guided by non-Markovian quantum dynamics can deliver a positive multifunctional net energy benefit (power $+$ water $+$ carbon $+$ diagnostics) in a field-constrained agrivoltaic stack, and which barriers remain physical rather than economic. Within the modeled \qty{500}{\m\squared} cooperative case, the answer is affirmative but conditional: dual-band filtering raises the forward transfer yield by \qty{25}{\percent} and extends coherence lifetime by \qty{50}{\percent}; confining NPoM operation to a \qty{1}{\percent} sentinel fraction preserves \qty{97.1}{\percent} global canopy yield while retaining SERS diagnostics; the FAO-56 smart shield cuts evapotranspiration by \qty{28}{\percent} with a Net Ecological Benefit of \qty{72.6}{\kg\of{\ce{CO2e}}\per\m\squared\per\yr}; and the cooperative case pays back in \qty{4.32}{\yr} at \qty{30}{\percent} subsidy. These claims are bounded by the $>\qty{90}{\percent}$ local NPoM transport penalty, the single-exciton and site-independent-bath approximations, the $\pm\qty{10}{\percent}$ microclimate uncertainty band, and the modeled (not field-validated) fouling and matrix-penalty assumptions. For deployment, the economic barrier is dominant and tractable: local-material construction and green-synthesis gold lower the capital cost from \qty{322}{\USD\per\m\squared} to \qty{140}{\USD\per\m\squared} (unsubsidized payback \qty{2.68}{\yr}). The physical barriers --- nanofabrication of sub-\qty{2}{\nano\meter} NPoM gaps, in-field fouling of the sentinel fraction, OPV encapsulation lifetime under tropical irradiance --- require process engineering rather than subsidies. Priority next steps are 2DES validation of the filter protocol, low-loss dielectric antenna alternatives, field fouling trials of the sentinel network, and a localized Au/graphene sensor supply chain for Sub-Saharan manufacturing.

% =============================================================================
% METHODS
% =============================================================================

\section{Methods}

\subsection{Proof-of-concept scenario}\label{sec:methods_poc}
A \qty{500}{\m\squared} high-tunnel floriculture module in Cameroon's Western Highlands (mean irradiance \qty{5.76}{\kWh\per\m\squared\per\day}) anchors all system-level accounts. Crop production is taken at the horticultural standard of \qty{8.2}{\kg\per\m\squared\per\yr}; no coherence-derived biomass gain is monetized. Full capital, OPEX, water-credit, and cooperative cash-flow tables appear in \Cref{SI-sec:socioeconomic}.

\subsection{Quantum dynamics}
The 8-site FMO Hamiltonian with non-Hermitian reaction-center trap and the $9\times 9$ NPoM-dressed Hamiltonian are specified in \Cref{SI-sec:hamiltonian}. Open-system dynamics are propagated with the PT-HOPS/SBD formalism (hierarchy depth $L=8$, Matsubara $K=2$, $\Delta t=\qty{0.2}{\fs}$, $T=\qty{295}{\kelvin}$, $\Gamma_{\mathrm{RC}}=\qty{0.15}{\per\fs}$). NPoM-on trajectories initialize the excitation in the plasmon mode while NPoM-off trajectories initialize on BChl site~1 (two illumination geometries); suppression values are conditional on this choice. Convergence, Hermiticity, positivity, and disorder statistics are reported in \Cref{SI-sec:convergence} and \Cref{SI-sec:comparative_hdf5}. Forward transfer yield is defined by the time-integrated trapping population on sites 3 and 4; coherence lifetime is the $1/e$ decay of the $l_1$-norm coherence.

\subsection{FAO-56 microclimate and lifecycle}
Crop reference and actual evapotranspiration are computed with the FAO-56 Penman--Monteith equation under the spectral transmittance of the OPV+NPoM shield. The resulting \qty{28}{\percent} ET reduction (central value) carries a $\pm\qty{10}{\percent}$ microclimate uncertainty band. Net Ecological Benefit aggregates avoided grid electricity, water savings, and embodied carbon of the module over a 20-year service life; full inventory and allocation rules are in the Supplementary Material.

\subsection{BB84 telemetry and SERS diagnostics}
A prepare-and-measure BB84 link secures the IoT telemetry channel; the fail-safe falls back to classical AES when the quantum bit-error rate exceeds the security threshold (\Cref{SI-sec:qkd}). SERS enhancement factors are obtained from the same NPoM mode-volume scan that supplies the transport yields; environmental robustness (temperature, humidity, soiling) is parameterized by the DynamicCalibrator module described in the SI.

\subsection*{Data availability and reproducibility}
All production trajectories, parameter files (v1.4), and analysis notebooks will be deposited in Zenodo with a DOI (HDF5 snapshot plus extraction script ready at submission). The digital-twin orchestrator and BB84 simulator are available from the corresponding author upon reasonable request.

\section*{Acknowledgements}
This work was supported by the University of Yaound\'{e}~I and the University of Douala.
Computational resources were provided by the penavoraserver cluster (ThinkStation P720, \qty{128}{\giga\byte} RAM, 48~cores).
We thank the MesoHOPS development team for numerical framework support.

\section*{CRediT authorship contribution statement}
\textbf{S.C. Teguia Kouam}: Conceptualization, Methodology, Software, Formal analysis, Writing -- original draft.
\textbf{T. Goumai Vedekoi}: Validation, Writing -- review \& editing.
\textbf{J.-P. Tchapet Njafa}: Validation, Writing -- review \& editing.
\textbf{J.-P. Nguenang}: Validation, Writing -- review \& editing.
\textbf{S.G. Nana Engo}: Supervision, Project administration, Writing -- review \& editing.

\section*{Declaration of competing interest}
The authors declare that they have no known competing financial interests or personal relationships that could have appeared to influence the work reported in this paper.

\section*{Funding}
This research received no specific grant from funding agencies in the public, commercial, or not-for-profit sectors. Institutional support was provided by the University of Yaound\'{e}~I and the University of Douala.

\section*{Declaration of generative AI and AI-assisted technologies in the writing process}
During the preparation of this work the authors used the following AI-assisted tools: OpenAI ChatGPT (GPT family large language model) and Anthropic Claude (Claude family large language model), applied solely to language editing, LaTeX formatting support, and reference-formatting checks. No scientific content, data, analysis, or interpretation was generated by these tools. All scientific content was reviewed and verified by the authors, who take full responsibility for the content of the publication.

\section*{Nomenclature}
\small
\begin{tabular}{@{}l p{0.68\linewidth}@{}}
\toprule
\textbf{Symbol} & \textbf{Meaning} \\
\midrule
$\Phi_{\mathrm{FT}}$ / $\Phi_{\mathrm{FT}}^{\mathrm{global}}$ & Local / area-weighted global forward transfer yield \\
$\tau_c$ & Exciton coherence lifetime ($1/e$ decay of $l_1$-norm coherence) \\
$g_0$ / $V_{\mathrm{mode}}$ & Single-emitter plasmon--exciton coupling / NPoM mode volume \\
$\Gamma_{\mathrm{RC}}$ & Reaction-center trapping rate \\
$\mathrm{EF}_{\mathrm{SERS}}$ & Relative Purcell-scaled SERS enhancement factor \\
$\mathrm{ET}_c$ / $K_c$ / $f_{\mathrm{shade}}$ & Crop evapotranspiration / crop coefficient / OPV shading factor \\
NEB & Net Ecological Benefit \\
QBER & Quantum bit error rate (BB84 link) \\
$\alpha$ & Sentinel area fraction \\
\bottomrule
\end{tabular}

% =============================================================================
% FIGURES
% =============================================================================

\begin{figure}[htbp]
\centering
\includegraphics[width=0.63\columnwidth]{Figure_Setup_Dispositif.png}
\caption{\textbf{Multi-scale quantum agrivoltaic digital twin.} Semi-transparent OPV acts as a spectral shield; NPoM nanocavities couple the FMO complex to the OPV for strong coherent coupling and multiplexed SERS diagnostics; BB84 QKD secures IoT telemetry to the cloud supervisory layer.}\label{fig:system_setup}
\end{figure}

\begin{figure}[htbp]
\centering
\includegraphics[width=\columnwidth]{Figure_Plateau_Canopy.png}
\caption{\textbf{Local quenching, global stability.} (a)~Forward transfer yield $\Phi_{\mathrm{FT}}$ across the NPoM mode-volume scan ($n=20$, $T=\qty{295}{\kelvin}$): local transport is quenched by more than \qty{90}{\percent} relative to the passive baseline ($0.980$), with a maximum of $0.0799$ at $V_{\mathrm{mode}}=\qty{1.2}{\nm\cubed}$ (\Cref{tab:npom_volume_scan}). (b)~Area-weighted global canopy yield for a \qty{1}{\percent} sentinel fraction (\Cref{eq:phi_global}): flat at $\num{0.971}$ across the scan --- a constant 0.9-point canopy-scale cost. The flat curve in (b) is the design-relevant outcome: diagnostics integrate without compromising device-level yield. Full provenance in \Cref{SI-sec:comparative_hdf5}.}
\label{fig:flat_canopy}
\end{figure}

\begin{figure}[htbp]
\centering
\includegraphics[width=\columnwidth]{Figure1_Quantum_Dynamics.png}
\caption{\textbf{Quantum dynamics of the 8-site FMO--NPoM stack.} (a)~Energy level diagram of the $9\times 9$ dressed Hamiltonian ($g_0$: plasmon--exciton coupling). (b)~Population dynamics at $V_{\mathrm{mode}}=\qty{1.2}{\nm\cubed}$, $T=\qty{295}{\kelvin}$ ($n=20$; excitation initialized in the plasmon mode, NPoM-off runs initialize on BChl site~1): the plasmon-like state retains \qty{83}{\percent} (final population \num{0.8256}) of the excitation at $t=\qty{1}{\ps}$, starving the trapping sites --- origin of the \qty{91.8}{\percent} local penalty. (c)~Cumulative $\Phi_{\mathrm{FT}}(t)$ saturating at \num{0.0799}. Production parameters $L=8$, $K=2$; full trajectories in \Cref{SI-sec:comparative_hdf5}.}
\label{fig:quantum_dynamics}
\end{figure}

\begin{figure}[htbp]
\centering
\includegraphics[width=\columnwidth]{Figure2_SERS_Readout.png}
\caption{\textbf{Dynamic photo-protection and diagnostics.} (a)~Floquet--Stark hysteretic switch: resonant channels close above \qty{850}{\W\per\m\squared} and reopen below \qty{750}{\W\per\m\squared} (solid-state analogue of non-photochemical quenching; design curves in \Cref{SI-sec:npq}). (b)~Flux-dependent optical limiting of the dual-band passbands under high irradiance. (c)~SERS enhancement vs.\ mode volume from the same NPoM scan that supplies the transport yields (\Cref{SI-sec:comparative_hdf5}).}\label{fig:npq_switch}
\end{figure}

\begin{figure}[htbp]
\centering
\includegraphics[width=\columnwidth]{Figure3_LCA_NEB_Comparison.png}
\caption{\textbf{Water--energy--food network assessment} for the modeled \qty{500}{\m\squared} cooperative module (Cameroon Western Highlands; \Cref{SI-sec:socioeconomic}). (a)~FAO-56 evapotranspiration: open field $\mathrm{ET}_c=\qty{4.51}{\mm\per\day}$ vs.\ OPV smart shield $\qty{3.25}{\mm\per\day}$ (\qty{28}{\percent} reduction). (b)~Net Ecological Benefit across three scenarios: Quantum OPV (A:~$\qty{72.6}{\kg\of{\ce{CO2e}}\per\m\squared\per\yr}$, biomass $\qty{13.0}{\kg\per\m\squared\per\yr}$ = \qty{12.0}{\kg\per\m\squared\per\yr} reference $\times$ capped yield factor \num{1.0} $\times$ \num{1.08} biostimulation boost, excluded from baseline revenue), static opaque shading (B:~higher carbon but suppressed canopy, \qty{10.0}{\kg\per\m\squared\per\yr} reference $\times$ \num{0.8} canopy factor), and open field (C, \qty{12.0}{\kg\per\m\squared\per\yr} reference $\times$ \num{0.7} canopy factor). Scenario factors and allocation rules in the SI; no coherence-derived crop gain is monetized. (c)~Undiscounted payback vs.\ blended-finance subsidy rate: \qty{4.32}{\yr} at \qty{30}{\percent} subsidy.}
\label{fig:lca_comparison}
\end{figure}

% =============================================================================
% REFERENCES
% =============================================================================

\bibliographystyle{elsarticle-num}
\bibliography{references}

\end{document}
